\section{部分子分布}

\subsection{一般矩阵元与部分子分布}

定义部分子分布的基本对象是一个双定域算符的矩阵元，
（略去其自旋结构中无关紧要的细节后）它可以一般地写成
$\langle  p | \phi (0)  \phi(z)  | p \rangle $。
由于 Lorentz 变换下的不变性，它是两个标量 $(pz)$
（记作 $-\nu$）与 $z^2$（或写成 $-z^2$，以便对类空的 \mbox{$z$ 取正值)}的函数：
  \begin{align}
  \langle p |   \phi(0) \phi (z)|p \rangle %\equiv B(z,p)
=  & {\cal M} (-(pz), -z^2)  =  {\cal M} (\nu, -z^2)
\,  .
 \
 \label{lorentz}
\end{align}
可以证明 \cite{Radyushkin:2016hsy,Radyushkin:1983wh}：对所有相关的 Feynman 图，
它对 $(pz)$ 的 Fourier 变换 ${\cal P} (x, -z^2)$ 具有支撑区 $-1 \leq x \leq 1$，即
   \begin{align}
 {\cal M} (-(pz), -z^2)
&   =
 \int_{-1}^1 dx
 \, e^{-i x (pz) } \,  {\cal P} (x, -z^2)  \   .
  \label{MPD}
\end{align}
式 (\ref{MPD}) 可作为 $x$ 的协变定义。
在这一 $x$ 的定义中，无需假设 $p^2=0$ 或 $z^2=0$。

选取类光的 $z$（例如只含光锥分量 $z_-$），我们用 twist-2 部分子分布
$f(x)$ 来参数化该矩阵元：
  \begin{align}
  {\cal M} (-p_+ z_- , 0)  =
   \int_{-1}^1 dx \, f(x) \,
e^{-ixp_+ z_-} \,  \   .
 \label{twist2par0}
\end{align}
此定义也可改写为
  \begin{align}
  {\cal M} (\nu, 0)  =
   \int_{-1}^1 dx \, f(x) \,
e^{ix\nu} \,  \  .
 \label{twist2par1}
\end{align}
逆关系为
\begin{align}
 f  (x)  =\frac{1}{2 \pi}  \int_{-\infty}^{\infty}  d\nu \, e^{-i x \nu}  \, {\cal M}(\nu , 0)    = {\cal P} (x, 0)
  \  .
\label{fxMnu}
\end{align}

由于 $f(x) =   {\cal P} (x, 0) $，函数 $ {\cal P} (x, -z^2) $ 把 PDF 的概念推广到了
非类光间隔 $z^2$ 的情形（原则上 $z^2$ 甚至可以是类时的）。
按照文献 \cite{Radyushkin:2017cyf}，我们称之为
{\it 伪部分子分布}（pseudo-PDF）。变量 $(pz)=-\nu $ 常被称为
{\it Ioffe 时间}\ \cite{Ioffe:1969kf}，
相应地 ${\cal M}(\nu , -z^2) $ 称为 {\it Ioffe 时间分布}\ \cite{Braun:1994jq}。


在可重正理论（包括 QCD）中，
函数 $ {\cal M} (\nu, -z^2)  $
具有 $\sim \ln (-z^2)$ 型对数奇异性，
正是这些奇异性生成了部分子密度的微扰演化。
在基于算符乘积展开（OPE）的方法中，标准做法是用某种方案
消除这些奇异性。
最流行的是基于维数正规化的 $\overline{\rm MS}$ 方案。
于是所得的 PDF 依赖于重正化标度 $\mu$，因此应把 PDF 写成 $f(x, \mu^2)$。

在小的类空 $z^2$ 处、领头对数精度下，伪部分子分布与 $\overline{\rm MS}$ 分布之间
仅差第二个宗量的简单重标度。
特别地，当 $z^2=-z_3^2$ 时有
\begin{align}
    {\cal P} (x, z_3^2) = f \left (x, (2 e^{-\gamma_E}/ z_3)^2 \right )
  \  ,
\label{Ptof}
\end{align}
其中 $\gamma_E$ 是 Euler 常数。
$\mu$ 与 $1/z_3$ 之间的重标度因子非常接近 1，因为 \mbox{$2e^{-\gamma_E}=1.12$}。


\subsection{横动量依赖分布与准分布}

把靶动量 $p$ 取为纵向的，\mbox{$p= (E, {\bf 0}_\perp, P)$}，
便可引入横向自由度。
特别地，取只含 $z_-$ 与 \mbox{$z_\perp= \{z_1,z_2\}$} 分量的 $z$，
可以定义 {\it TMD}\ ${\cal F} (x, k_\perp^2)$：
      \begin{align}
 {\cal P} (x,  z_\perp^2)
&   =
    \int  d^2{\bf k}_\perp   e^{i ({\bf k}_\perp {\bf z}_\perp)}
      {\cal F} (x, {k}_\perp^2)
 \   .
  \label{MTMD0}
\end{align}
在此语境下，伪部分子分布 $ {\cal P} (x,  z_\perp^2) $ 实际上就是许多 TMD 研究中常见的
{\it 撞击参数分布}（impact parameter distributions）。


由于格点上无法安排类光间隔，文献 \cite{Ji:2013dva} 建议考虑等时的类空间隔
$z= (0,0,0,z_3)$（简称 \mbox{$z=z_3$}）。于是
在 \mbox{$p=(E, 0_\perp, P)$} 系中，
可以通过如下参数化引入准部分子分布 $Q(y, P)$：
  \begin{align}
  \langle p |   \phi(0) \phi (z_3)|p \rangle %\equiv B(z,p)
=  &
\int_{-\infty}^{\infty}   dy \,
 Q(y, P) \,  e^{i y  P z_3 } \,
 \  .
 \label{newVDFxzQ}
\end{align}
按此定义，准部分子分布 $Q(y,P)$ 描述的是部分子携带母强子第三动量分量 $P$
的份额 $y$ 的概率。
此时变量 $\nu$ 与 $-z^2$ 由 $Pz_3$ 与 $z_3^2$ 给出，因此
 \begin{align}
{\cal M} (\nu,z_3^2)
=  &
\int_{-\infty}^{\infty}   dy \,
 Q(y, P) \,  e^{i y  \nu } \,
 \  .
\end{align}
由于 $z_3^2= \nu^2/P^2$，逆 Fourier 变换可写为
\begin{align}
  Q(y,  P)   =\frac{1}{2 \pi}  \int_{-\infty}^{\infty}  d\nu \,
   \, e^{-i y  \nu}  \, {\cal M} (\nu,  \nu^2/P^2)    \  .
\label{IxM}
\end{align}
它表明当 $P \to \infty$ 时 $Q(y,  P) $ 趋于 $f(y)$，
因为形式上当 \mbox{$P \to \infty$} 时
${\cal M} (\nu,  \nu^2/P^2) \to {\cal M} (\nu, 0)$。

如文献 \cite{Radyushkin:2016hsy} 所确立的，准部分子分布可用 TMD 写出：
 \begin{align}
 Q(y, P) /P =  & \,\int_{-\infty}^{\infty} d  k_1
\int_{-1} ^ {1} d  x \,   {\cal F} (x, k_1^2+(y-x)^2P^2 )
  \  .
 \label{QTMD}
 \end{align}

\subsection{量子色动力学（QCD）情形}

对于 QCD 的非单态部分子密度，
所考虑的矩阵元为
    \begin{align}
 {\cal M}^\alpha  (z,p) \equiv \langle  p |  \bar \psi (0) \,
 \gamma^\alpha \,  { \hat E} (0,z; A) \psi (z) | p \rangle \  ,
\label{Malpha}
\end{align}
其中 $ { \hat E}(0,z; A)$ 是夸克（基础）表示中标准的 $0\to z$ 直线规范链接。
这些矩阵元可分解为 $p^\alpha$ 与 $z^\alpha$ 两部分：
\begin{align}
{\cal M}^\alpha  (z,p) = &2 p^\alpha  {\cal M}_p (-(zp), -z^2)
 + z^\alpha  {\cal M}_z (-(zp),-z^2)
\ .
\end{align}
只有 ${\cal M}_p (-(zp), -z^2) $ 部分在 $z^2 \to 0$ 时给出 twist-2 分布。

引入 TMD 时，取 $z=(z_-, z_\perp)$ 以及 ${\cal M}^\alpha$ 的 \mbox{$\alpha=+$} 分量，
此时 $z^\alpha$ 部分自动消失。
此后 $ {\cal M}_p(\nu, z_\perp^2)$ 成为 ${\cal M}^\alpha  (z,p)$ 中唯一存留的部分，
在余下的讨论中我们采用简写 ${\cal M} \equiv {\cal M}_p$。

对于准分布 $Q(y,P)$，
可以考虑 \mbox{$ {\cal M}^\alpha  (z=z_3,p)$} 的时间分量来避开 $z^\alpha$ 污染，
并定义
 \begin{align}
&  {\cal M}^0   (z_3,p)
  =  2p^0
 \int_{-1}^1 dy\,
Q(y,P) \,
 \,  e^{i  y Pz_3 }  \  .
 \label{OPhixspin12}
\end{align}



\subsection{因子化模型}

准部分子分布的结构可以用最简单的模型来说明：
其中 TMD ${\cal F}  (x, k_\perp^2)$ 的非微扰（软）部分表示为一个乘积
    \begin{align}
{\cal F}^{\rm soft}  (x, k_\perp^2) = f(x) K(k_\perp^2)
\end{align}
——共线部分子分布 $f(x)$ 与通常取高斯型的 \mbox{$k_\perp^2$ 依赖}因子
$K(k_\perp^2)$ 的乘积。
我们将看到，即便由这些简单的因子化模型出发，
准部分子分布也会呈现出相当复杂的结构。

对 Ioffe 时间分布 ${\cal M} (\nu,-z^2)$ 而言，该假定对应于其软部分的
因子化假设
   \begin{align}
   {\cal M}^{\rm soft}  (\nu,z_3^2) = {\cal M}^{\rm soft}  (\nu,0){\cal M} (0,z_3^2) \
    \end{align}
即便软的 TMD 因子化，准部分子分布的软部分仍具有式 (\ref{QTMD}) 的卷积结构。
例如，取高斯形式
  \begin{align}
K_G(k_\perp^2) =  &
\frac{1}{\pi  \Lambda^2}  e^{-k_\perp^2/\Lambda^2}
  \  ,
 \label{DTMDsmallfac}
\end{align}
便得到如下准部分子分布模型
 \begin{align}
 Q_G(y, P)  = &\frac{P}{\Lambda \sqrt{\pi} }  \,
 \int_{-1}^1 dx\,  %
f(x) \,
  e^{- (x -y)^2 P^2 / \Lambda^2 }
 \  .
 \label{QinG}
\end{align}
取一个形似核子价密度的简单玩具 PDF
$f(x)=4(1-x)^3 \theta (0\leq x \leq 1)$ 作为 $f(x)$，
就得到图 \ref{Qgy} 所示的曲线。
大 $P$ 时准部分子分布明显趋于 $f(y)$ 的 PDF 形状。
但只有当 $P \sim 10 \Lambda$ 时，
得到的准部分子分布才比较接近 $P \to \infty$ 的极限形状。
鉴于 $\Lambda \sim \langle k_\perp \rangle$，
\mbox{$P\sim 10 \Lambda$} 的估计换算过来就是 $P \sim 3$ GeV，
这是相当大的动量。


 \begin{figure}[t]
    \centerline{\includegraphics[width=3.6in]{images/QyPn.pdf}}
    \caption{因子化高斯模型中准部分子分布 $Q(y,P)$ 的演化，对应 {$P/\Lambda =1, 5, 10, 50$}。
    \label{Qgy}}
    \end{figure}
